\documentclass[sigconf,nonacm]{acmart}
\settopmatter{printacmref=false}
\usepackage{subcaption}
\usepackage{booktabs}
\usepackage{float}
\usepackage{placeins}
\renewcommand{\topfraction}{0.9}
\renewcommand{\bottomfraction}{0.8}
\setcounter{topnumber}{2}
\setcounter{bottomnumber}{2}
\setcounter{totalnumber}{4}
\setcounter{dbltopnumber}{2}
\renewcommand{\dbltopfraction}{0.9}
\renewcommand{\textfraction}{0.07}

%%
%% \BibTeX command to typeset BibTeX logo in the docs
\AtBeginDocument{%
  \providecommand\BibTeX{{%
    Bib\TeX}}}

%% Rights management information.  This information is sent to you
%% when you complete the rights form.  These commands have SAMPLE
%% values in them; it is your responsibility as an author to replace
%% the commands and values with those provided to you when you
%% complete the rights form.
\setcopyright{none}
\copyrightyear{2025}
\acmYear{2025}
\acmDOI{}

%% These commands are for a PROCEEDINGS abstract or paper.
\acmConference[15-721]{Advanced Database Systems}{Spring 2025}{Carnegie Mellon University}

\acmISBN{}


\begin{document}

\title{Parallel Bloom Join Implementation Summary}

\author{Kartica Modi}
\affiliation{%
  \institution{Carnegie Mellon University}
  \city{Pittsburgh}
  \state{PA}
  \country{USA}
}

\author{Aarrya Saraf}
\affiliation{%
  \institution{Carnegie Mellon University}
  \city{Pittsburgh}
  \state{PA}
  \country{USA}
}

\renewcommand{\shortauthors}{Modi and Saraf}

\begin{abstract}
The Parallel Bloom Join optimization aims to reduce the volume of data processed during a Hash Join by filtering out non-matching tuples from the outer relation early in the pipeline (at the scan level). Using a shared Bloom filter in parallel execution, we observe execution time reductions of up to 40\% in low-selectivity workloads by reducing I/O and tuple transfer overhead. We further introduce adaptive gating to dynamically disable the filter when ineffective.
\end{abstract}

%%
%% The code below is generated by the tool at http://dl.acm.org/ccs.cfm.
%% Please copy and paste the code instead of the example below.
%%
\begin{CCSXML}
<ccs2012>
<concept>
<concept_id>10002951.10002952.10003190</concept_id>
<concept_desc>Information systems~Database management system engines</concept_desc>
<concept_significance>500</concept_significance>
</concept>
<concept>
<concept_id>10002951.10002952.10003190.10003192</concept_id>
<concept_desc>Information systems~Database query processing</concept_desc>
<concept_significance>500</concept_significance>
</concept>
</ccs2012>
\end{CCSXML}

\ccsdesc[500]{Information systems~Database management system engines}
\ccsdesc[500]{Information systems~Database query processing}

\keywords{PostgreSQL, Bloom Filter, Hash Join, Query Optimization, Parallel Execution, Adaptive Processing}

\maketitle

\section{Introduction}
Analytic queries in modern database systems often involve large-scale joins between fact and dimension tables. The Hash Join algorithm is the standard workhorse for such equality joins, where an in-memory hash table is built from the smaller (inner) relation, and the larger (outer) relation is scanned to probe for matches. However, when the outer relation is massive, the overhead of scanning non-matching rows becomes a significant bottleneck. This cost is not merely I/O; in PostgreSQL, simply retrieving a tuple from a heap page, projecting it into a \texttt{TupleTableSlot}, and deforming its attributes incurs substantial CPU overhead (interpretation and parsing). For joins with low selectivity (e.g., 5\%), the system wastes 95\% of its cycles deforming tuples that will eventually be discarded by the join operator.

To mitigate this, we propose a \textbf{Parallel Bloom Join} optimization for PostgreSQL. By constructing a Bloom filter on the inner relation's keys and pushing this filter down to the sequential scan of the outer relation, we can discard non-matching tuples early in the pipeline. Note that our current implementation filters tuples immediately after the sequential scan retrieves them (avoiding inter-operator communication), but it does not bypass the initial \texttt{TupleTableSlot} projection.

Our contributions are as follows:
\begin{enumerate}
    \item \textbf{Scan-Level Pruning}: Moving the filtering logic from the join operator to the scan operator (`SeqScan`), minimizing tuple transport cost.
    \item \textbf{Parallel-Aware Design}: Implementing both shared and independent Bloom filter strategies to support PostgreSQL's parallel query execution.
    \item \textbf{Adaptive Gating}: Dynamically disabling the filter if logical selectivity is found to be low, avoiding regression overhead.
\end{enumerate}

\paragraph{Background: Hash Joins in PostgreSQL}
The standard PostgreSQL Hash Join implementation proceeds in distinct phases. In the simplest in-memory case, the executor scans the inner relation to build a hash table, then scans the outer relation to probe it. When the inner relation size exceeds \texttt{work\_mem}, the system switches to a Hybrid Hash Join strategy, partitioning both relations into batches that are spilled to disk and joined independently. In parallel execution, multiple workers cooperate to build a shared hash table (or shared batch files) in dynamic shared memory, synchronizing via barriers before proceeding to the probe phase. Our work builds directly upon this parallel infrastructure.

\section{Related Work}
The application of probabilistic data structures to optimize database operations is a foundational concept. Bloom \cite{bloom1970space} introduced the Bloom filter as a space-efficient structure for membership testing. In the domain of distributed database systems, Mackert and Lohman \cite{mackert1986r_optimizer} leveraged this structure in the $R^*$ optimizer to define the ``Bloomjoin,'' which minimizes network transmission by filtering tuples at the source before shipment.

Our approach aligns with the broader paradigm of Adaptive Query Processing (AQP) and Sideways Information Passing (SIP), where information gathered during query execution is used to optimize the remaining workload. Ives et al. \cite{ives2009adaptive} demonstrated the efficacy of passing filters between operators to prune irrelevant data streams. Similarly, Zeller and Gray \cite{zeller1990adaptive} introduced adaptive hash join algorithms to robustly handle varying memory constraints, a precursor to the Hybrid Hash Join employed by PostgreSQL.

Within the PostgreSQL ecosystem, Vondra \cite{vondra2015bloomhashbatches} proposed a preliminary implementation of Bloom filtering for batched hash joins. However, this work was limited to single-threaded execution and performed filtering within the join operator itself. Crucially, the patch was ultimately rejected because it relied on optimizer estimates for selectivity, which are notoriously prone to error. In cases where the optimizer underestimated selectivity, the overhead of building and checking the filter caused performance regressions, a risk the community deemed unacceptable. Our work bridges this gap by implementing a parallel-aware, shared Bloom filter mechanism and proposing robust adaptive safeguards.

\section{Implementation}

We integrate a Bloom filter into the query execution pipeline to prune tuples as early as possible. This involves extensions to the Bloom filter data structure, the Hash Join executor, and the Sequential Scan logic.

\subsection{Shared Bloom Filter Architecture}
To support parallel query execution, we enhanced the standard Bloom filter to reside in memory allocated via the \textbf{Dynamic Shared Area (DSA)} allocator, backed by dynamic shared memory (\texttt{DSM}) segments. The filter is initialized using atomic operations (specifically \\ \texttt{pg\_atomic\_fetch\_or\_u32}) to ensure thread-safe bit setting during the build phase. We utilize $k=3$ hash functions by default, balancing CPU cost against collision rates.

Parallel workers cooperatively populate the filter while building the inner hash table. This design ensures all workers scanning the outer table benefit from a unified global filter state. However, this shared write access introduces significant cache coherence overhead. This manifests as both \textbf{True Sharing} (multiple workers writing to the same word) and \textbf{False Sharing} (writing to different words within the same 64-byte cache line). In both cases, the hardware cache coherence protocol (MESI) triggers expensive invalidations and bus traffic, degrading performance despite the lack of logical lock contention.

\subsection{Pipeline Integration}
The optimization creates a feedback loop between the consumer (Hash Join) and producer (SeqScan) operators.
\paragraph{Build Phase}
During the hash table build, keys are concurrently inserted into the Bloom filter. Once the build phase concludes, the fully populated filter is "attached" to the outer relation's scan node.

\paragraph{Probe \& Prune}
The Sequential Scan operator incorporates a probe step. Currently, this check is performed immediately after \texttt{table\_scan\_getnextslot()} retrieves and materializes a tuple into a slot.
\begin{itemize}
    \item \textbf{Hit}: The tuple is passed up to the join operator for full verification.
    \item \textbf{Miss}: The tuple is immediately discarded. While this does not skip the scan-level slot projection, it prevents the tuple from entering the upper executor nodes, saving significant transport and qualification overheads.
\end{itemize}

\subsection{Adaptive Gating Mechanism}
Bloom filters introduce a small CPU overhead per tuple. If the filter selectivity is poor (i.e., most tuples pass), this overhead contributes to regression. We implement an adaptive gating mechanism: if the first 100 probed tuples all pass the filter, the system infers low pruning power and disables the filter for the remainder of the scan, reverting to standard execution.

\subsection{Evolution of Parallel Strategy}
Our initial design utilized a single shared Bloom filter protected by barriers. While correct, this introduced significant latch contention. We are currently transitioning to a "per-worker" independent filter strategy to eliminate synchronization overheads, trading slightly higher memory usage for linear scaling.




\section{Evaluation}

We evaluate the efficacy of the Parallel Bloom Join through two distinct studies: a baseline replication and a controlled ablation study.

\subsection{Methodology}
Experiments were conducted on a standard PostgreSQL development environment hosted on AWS. We used an \textbf{m6id.4xlarge} instance to ensure sufficient compute and I/O throughput.
\begin{itemize}
    \item \textbf{Processor}: 16 vCPUs (Intel Xeon Scalable, x86\_64).
    \item \textbf{Memory}: 64 GiB.
    \item \textbf{Storage}: 950 GB NVMe SSD (Instance Store) for data, 8 GB gp3 for root.
    \item \textbf{OS}: Ubuntu Server 24.04 LTS.
\end{itemize}
PostgreSQL was configured with a \texttt{work\_mem} of 64MB to forcefully trigger spill-to-disk behavior for our larger dataset tests, ensuring we evaluated both memory-resident and I/O-bound regimes.

We define \textit{Selectivity} as the fraction of outer relation tuples that typically match a tuple in the inner relation.

\subsubsection{Baseline Replication}
To validate our system against historical attempts, we replicated the methodology of Vondra. We generate a synthetic dataset with 10M rows, duplicated 10x for the outer table (100M rows). Joins are performed on the integer row indices. This workload characterizes CPU-bound join performance. 
Crucially, we artificially vary selectivity in this test case by filtering on the MD5 hash itself. We filter the inner table to only include rows where the first character of the MD5 hash is less than a specific hex parameter. This method, employed by Vondra, changes both the effective join selectivity and the size of the inner table simultaneously.

\subsubsection{Controlled Ablation Study}
We developed a custom benchmark harness (\texttt{benchmark.py}) to isolate variables.
\paragraph{Data Generation}
We utilize synthetic Fact and Dimension tables where \texttt{int\_value} serves as the join key. The Fact table generation logic guarantees precise selectivity (e.g., 20\% matches) by drawing non-matching keys from a disjoint range logic.

\paragraph{Dynamic Configuration}
To facilitate extensive ablation without the overhead of recompiling the database server, we exposed several custom configuration parameters (GUCs) within PostgreSQL. These settings allowed us to dynamically tune critical system behaviors at runtime:
\begin{enumerate}
    \item \textbf{Bloom Filter Parameters}: We added controls to adjust the \texttt{bloom\_filter\_multiplier} and the number of hash functions ($k$) on a per-session basis.
    \item \textbf{Strategy Toggles}: Boolean flags were introduced to enable or disable the Bloom filter entirely, as well as to toggle between the Shared and Independent Bloom Filter implementations.
\end{enumerate}
This flexibility was instrumental in generating the wide array of performance profiles presented in our results.

\paragraph{Decision Variables}
We isolate the performance impact of five key decision variables:
\begin{itemize}
    \item \textbf{Inner Table Size}: The size of the build relation.
    \item \textbf{Outer Table Size}: The size of the larger probe relation.
    \item \textbf{Bloom Filter Size}: Controlled via the \\ \texttt{bloom\_filter\_multiplier} parameter. This value is defined as the Postgres estimated number of unique keys in the inner relation multiplied by this multiplier.
    \item \textbf{Hash Functions}: The number of hash functions ($k$) used for the filter.
    \item \textbf{Selectivity}: The fraction of outer table tuples that find a match in the inner table.
\end{itemize}

\subsection{Results}
In this section, we present the core performance analysis of our Parallel Bloom Join implementation, focusing on single-threaded efficiency and parallel scalability limits.

\paragraph{Single-Threaded Performance}
In serial execution (Figure~\ref{no_parallel}), the Parallel Bloom Join demonstrates substantial gains. We observe large improvements in execution time compared to the baseline, particularly at lower selectivities where the filter effectively prunes non-matching tuples.

\begin{figure}[ht]
    \centering
    \includegraphics[width=\linewidth]{15721_final_report/figures/no_parallel.png}
    \caption{Bloom Join on Vondra's Test Case; Single Processed}
    \label{no_parallel}
    \Description{Visualization of Bloom Join on Vondra's Test Case with a single process.}
\end{figure}

\paragraph{Parallel Scalability}
The scalability results reveal a diminishing return on parallelism. With 2 workers (Figure~\ref{Vondra, 2 parallel}), we observe some improvement over the baseline. However, as we scale to 4 workers (Figure~\ref{Vondra, 4 parallel}), the gains become marginal. This suggests that while the shared Bloom filter provides benefits, the coordination overhead or lock contention may be limiting scalability at higher concurrencies.

\begin{figure}[ht]
    \centering
    \includegraphics[width=\linewidth]{15721_final_report/figures/2paralle.png}
    \caption{Bloom Join on Vondra's Test Case; 2 Parallel Workers}
    \label{Vondra, 2 parallel}
    \Description{Bloom Join on Vondra's Test Case; 2 Parallel Workers}
\end{figure}

\begin{figure}[ht]
    \centering
    \includegraphics[width=\linewidth]{15721_final_report/figures/4parallel.png}
    \caption{Bloom Join on Vondra's Test Case; 4 Parallel Workers}
    \label{Vondra, 4 parallel}
    \Description{Bloom Join on Vondra's Test Case; 4 Parallel Workers}
\end{figure}

\subsection{Comparison with Historical Data}
For context, we present Vondra's original results in Figure~\ref{fig:vondra_res}. While hardware may differ, the trend confirms that bloom filtering offers significant gains when selectivity allows for effective pruning. Our results validate that this benefit persists in modern PostgreSQL versions, though we identify new bottlenecks in the parallel regime that were not present in Vondra's single-threaded implementation.

\begin{figure}[ht]
    \centering
    \includegraphics[width=\linewidth]{15721_final_report/figures/vondra_res.jpeg}
    \caption{Vondra's Original Performance Graph}
    \label{fig:vondra_res}
    \Description{A line graph comparing execution time between standard Hash Join and Bloom Join for varying selectivities (6.25\% to 37.5\%). The Bloom Join consistently outperforms the standard join.}
\end{figure}

\subsection{Ablation Study}
We further investigate the impact of the Bloom filter size on performance, controlled by the \texttt{bloom\_filter\_multiplier}. This parameter dictates the allocated memory for the filter relative to the estimated number of distinct keys in the inner relation. Figure~\ref{fig:bf_mult} illustrates the execution time across varying inner and outer table sizes.

\paragraph{Filter Size Trade-off}
The results indicate a "sweet spot" for the multiplier between 1.0x and 1.1x. This can be attributed to two countervailing factors:
\begin{enumerate}
    \item \textbf{False Positive Rate (FPR)}: Smaller multipliers (< 1.0) result in a denser bit array, increasing the collision rate. A high FPR allows more non-matching tuples to pass through the filter, reducing the pruning benefit and increasing downstream processing costs.
    \item \textbf{Cache Locality}: Larger multipliers (> 1.1) produce a sparser but larger bit array. As the filter size grows, it may exceed the capacity of the L3 cache. This incurs expensive cache misses during probing, increasing the CPU cost per tuple and negating the savings from pruning.
\end{enumerate}

\begin{figure}[ht]
    \centering
    \includegraphics[width=\linewidth]{15721_final_report/figures/bf_mult.png}
    \caption{Impact of Bloom Filter Multiplier on Execution Time}
    \label{fig:bf_mult}
    \Description{Graph showing execution time vs. multiplier. The curve is U-shaped with a minimum around 1.0-1.1. It shows instability when table sizes match.}
\end{figure}

\paragraph{Equal-Size Anomaly}
We observe erratic performance behavior when the inner and outer table sizes are comparable. In this regime, the Bloom filter is sufficiently large (proportional to the large inner table) to generally exceed the L3 cache size. This causes significant Last Level Cache (LLC) thrashing: the randomly accessing probe phase of the Bloom filter competes for cache lines with the hash table probes. Consequently, the memory bandwidth becomes saturated by these dual random-access patterns, and the pre-filtering cost outweighs the savings from skipping the join probe, as the CPU stalls on memory fetches.

\paragraph{Hash Function Sensitivity}
We further explore the interaction between filter size and the number of hash functions ($k$) in Figure~\ref{fig:bf_param_mat}. The data suggests that a multiplier of 1.1x combined with $k=2$ hash functions yields the optimal performance for larger joins. However, contrary to theoretical expectations, using a single hash function ($k=1$) proves surprisingly effective, particularly for smaller join sizes (as seen in Figure~\ref{fig:bf_mult}). In these scenarios, the CPU overhead of computing additional hashes and probing more bits outweighs the marginal gain in pruning power. This implies that "undershooting" the number of hash functions is generally a safer strategy than overshooting, as the penalty for a slightly higher false positive rate is less severe than the computational cost of complex hashing. Ideally, the system should \textbf{adaptively tune} $k$ based on the estimated cardinality and system load.

\begin{figure}[ht]
    \centering
    \includegraphics[width=\linewidth]{15721_final_report/figures/bf_param_mat.png}
    \caption{Hash Function vs. Multiplier Sensitivity}
    \label{fig:bf_param_mat}
    \Description{Heatmap showing execution time for combinations of multiplier sizes and hash function counts. The sweet spot is at 1.1x multiplier and 2 hash functions.}
\end{figure}

\subsection{Scalability and Contention}
We conclude our evaluation by analyzing the scalability limits of the shared Bloom filter architecture in Figure~\ref{fig:workers_plot}.
\begin{figure*}[ht]
    \centering
    \includegraphics[width=0.8\linewidth]{15721_final_report/figures/hash_join_workers_plot.png}
    \caption{Impact of Worker Count and Selectivity on Execution Time}
    \label{fig:workers_plot}
    \Description{Graph comparing execution time for 0, 2, and 4 workers across varying selectivities. 0 workers shows improvement; 2 and 4 workers show regression or stagnation due to contention.}
\end{figure*}

The single-threaded case (0 workers) exhibits the expected behavior: execution time is significantly improved at low selectivities and converges towards the baseline as selectivity approaches 0.8, where the filter's pruning power diminishes. However, with 2 and 4 parallel workers, we observe negligible performance gains compared to the baseline, and in some cases, degradation.

This stagnation is directly attributed to lock contention and coherence traffic on the shared Bloom filter. As workers concurrently attempt to atomic set and check bits, the resulting memory bus saturation from MESI protocol invalidations negates the potential speedup from parallelism. This result empirically validates our ongoing architectural shift towards Independent (Per-Worker) Bloom Filters to eliminate this coordination bottleneck.

\subsection{Investigation: Independent Filter Strategies}
Motivated by the scalability limits of the shared design, we implemented an experimental prototype using \textbf{Independent (Per-Worker) Bloom Filters}. In this model, each worker populates a local filter during the build phase, eliminating synchronization. These local filters are then merged into a global filter for the scan phase. We evaluated two merge strategies:

\paragraph{Atomic Merge (32-bit OR)}
Workers concurrently merge their local filters into the global, shared filter using \texttt{atomic\_or} operations on 32-bit words. As shown in Figure~\ref{fig:indep_bf}, this approach yielded results \textbf{very similar to just the shared bloom filter}, with negligible performance gains. The bottleneck essentially shifted from the Build Phase to the Merge Phase. The "Write Storm" during merging caused massive cache line thrashing, as multiple workers simultaneously attempted to update the same memory addresses in the shared filter.

\begin{figure}[ht]
    \centering
    \includegraphics[width=\linewidth]{15721_final_report/figures/indep_bf.jpeg}
    \caption{Independent Filters with Atomic Merge (Write Storm)}
    \label{fig:indep_bf}
    \Description{Graph showing poor performance of Atomic Merge strategy due to cache thrashing.}
\end{figure}

\paragraph{Locked Merge (Serialized)}
To mitigate thrashing, we serialized the merge phase by acquiring a lock on the global filter (Figure~\ref{fig:indep_bf_locked}). This successfully stabilized performance and \textbf{saw some improvements} over the atomic strategy. However, the serial nature of this merge process introduces an additional overhead component. While the build phase benefits from parallelism, the time required to sequentially merge large local filters eventually limits the overall speedup, preventing ideal linear scalability.

\textbf{Crucially, largely due to this merge overhead, both independent filter prototypes performed worse than the standard execution without any Bloom filters.} The cost of managing and merging these structures in parallel outweighed the pruning benefits in this high-contention regime.

\begin{figure}[ht]
    \centering
    \includegraphics[width=\linewidth]{15721_final_report/figures/indep_bf_locked.jpeg}
    \caption{Independent Filters with Locked Merge (Serialized)}
    \label{fig:indep_bf_locked}
    \Description{Graph showing limited performance using Locked Merge due to serialization overhead.}
\end{figure}

\FloatBarrier
\FloatBarrier
\section{Future Work}
Our evaluation highlights clear avenues for architectural refinement. While the current shared-filter implementation proves the viability of scan-level pruning, next-generation optimizations must address the scalability bottlenecks we identified.

\subsection{Optimized Merge Strategies}
The failure of our naive "Independent Filter" experiments (Atomic and Locked Merge) highlights that lock contention simply shifted from the Build Phase to the Merge Phase. To truly unlock scalability, we propose a \textbf{Leader Merge after Barrier} strategy.

In this design, each worker populates a strictly local Bloom filter during the build phase, eliminating all shared memory traffic. At the conclusion of the build phase, workers synchronize at a parallel barrier. A designated leader (e.g., the first worker) then performs the merge:
\begin{enumerate}
    \item The leader clears the global filter memory.
    \item It iterates through each worker's local filter, performing a bitwise OR into the global filter.
    \item All other workers wait at the barrier until the global filter is finalized.
\end{enumerate}

This approach eliminates the need for locks and atomic operations entirely during the merge. The merge process becomes a purely single-threaded, cache-friendly streaming operation. While the computational cost is linear ($O(n_{workers} \times n_{words})$), for typical filter sizes (megabytes), this duration is negligible compared to the total join time. This strategy prioritizes predictability and cache locality over theoretical parallelism in the brief merge phase.

\subsection{SIMD and Blocked Bloom Filters}
The current implementation relies on scalar execution for hashing and probing. Modern CPUs offer wide SIMD registers (e.g., AVX-512) that can significantly accelerate these operations. We plan to implement \textbf{Blocked Bloom Filters}, where the bit array is partitioned into cache-line-sized blocks.
\begin{itemize}
    \item \textbf{Vectorized Hashing}: Using SIMD instructions to compute $k$ hash values in parallel.
    \item \textbf{Cache-Local Probing}: Ensuring that all $k$ bits for a given key reside within a single cache line (block), reducing the probe cost to exactly one cache miss per tuple.
\end{itemize}

\subsection{Spill-File Pruning}
We propose extending our optimization to the Hybrid Hash Join scenario by implementing Bloom filters for each spill file. Constructing and writing this filter to disk incurs a linear cost of $1 \dots N$ I/Os. However, during the join phase, the filter is read into memory in chunks. For a buffer size of $B$, we pay 1 I/O to read the filter block. Since pruning a single tuple prevents it from being written to and subsequently read from a spill file, it saves exactly 2 I/Os. Thus, if the filter allows us to prune even one tuple per buffer (on average), the optimization breaks even; any further pruning results in a net performance gain. Furthermore, because these filters reside on disk and are streamed, they do not contend for main memory, alleviating the pressure observed in our in-memory experiments.

\subsection{Historical Selectivity Tracking}
To address the stability concerns that stalled previous efforts, we propose a persistent mechanism for learning join selectivities. We envision a global system catalog table maintaining the historical selectivity observed between every pair of joined tables.
\begin{enumerate}
    \item \textbf{Cold Start}: If no entry exists for a pair (selectivity is \texttt{NULL}), the system aggressively disables the Bloom filter to avoid regression, proceeds with a standard standard join, and records the observed selectivity.
    \item \textbf{Adaptive Execution}: If an entry exists, we check a threshold (e.g., 0.6). If the historical selectivity is $< 0.6$, we enable the Parallel Bloom Join. If $\ge 0.6$, we disable it.
    \item \textbf{Continuous Learning}: In all cases, the actual selectivity observed during execution is used to update the historical record (e.g., via an exponentially weighted moving average).
\end{enumerate}
This feedback loop ensures the system adapts to data updates over time and only pays the overhead of filtering when history suggests a high probability of benefit.

\section{Conclusion}
This work presents a comprehensive implementation and evaluation of a Parallel Bloom Join for PostgreSQL. By pushing membership testing down to the sequential scan operator, we successfully reduce the volume of data flowing through the execution pipeline.

Our experimental results demonstrate:
\begin{enumerate}
    \item \textbf{Efficacy}: In single-threaded execution, the optimization delivers substantial speedups, validating the theoretical benefit of scan-level pruning.
    \item \textbf{Parameter Sensitivity}: We established a rigorous tuning profile, identifying a Bloom Filter multiplier of \textbf{1.1x} combined with \textbf{$k=2$} hash functions as the robust "sweet spot" for maximizing performance while minimizing cache pressure (Figure~\ref{fig:bf_param_mat}).
    \item \textbf{Scalability Limits}: Our scalability analysis revealed that a shared-everything architecture suffers from atomic contention beyond 2 workers. This negative result is a critical finding, directing future engineering efforts toward shared-nothing (independent) filter designs.
\end{enumerate}

Ultimately, the Parallel Bloom Join represents a high-impact optimization for analytic workloads. With the proposed transition to independent, vectorized filters, this technique has the potential to become a standard feature in the PostgreSQL query engine.

\bibliographystyle{ACM-Reference-Format}
\bibliography{references}

\appendix
\input{lip_appendix}

\end{document}
